\origpage{261}{GRAPH MINORS. VIII}
\textit{Proof.} If $\mathscr{A}$ has no index, then there is a countable sequence $(G_1, T_1)$, $(G_2, T_2), \dots$ of members of $\mathscr{A}$ such that $|T_1| < |T_2| < \dots$. But then for $1 \le i < i'$, $(G_i, T_i)$ is not simulated in $(G_{i'}, T_{i'})$, because (S6) is not satisfied. Hence $\mathscr{A}$ is not well-rooted. $\square$

We mention that we believe that the converse to (4.1) also holds, and hope to publish a proof of this in a later paper. This would clearly imply Wagner's conjecture.

\stmt
(4.2) \textit{If $\mathscr{A}_1, \mathscr{A}_2$ are well-rooted then so is $\mathscr{A}_1 \cup \mathscr{A}_2$.}

\textit{Proof.} Any countable sequence of elements of $\mathscr{A}_1 \cup \mathscr{A}_2$ has a countable subsequence either with all entries in $\mathscr{A}_1$ or with all entries in $\mathscr{A}_2$. The result follows. $\square$

Grafts are one way to work with multiply-rooted graphs. But sometimes we need to distinguish the roots, and it is more convenient to work with $k$-tuples of vertices rather than subsets. A rooted graph then is a pair $(G, \tau)$, where $G$ is a graph and $\tau$ is a finite sequence of vertices of $G$. The index of a rooted graph is the length of $\tau$. If $\tau$ has length $k$ and $1 \le i \le k$, $\tau(i)$ denotes the $i$th term of $\tau$. If $(G, \tau)$ and $(H, v)$ are rooted graphs, a simulation of $(H, v)$ in $(G, \tau)$ is a simulation $\sigma$ of $H$ in $G$ with the following additional property:

\hspace*{12pt}(S7) $(G, \tau)$ and $(H, v)$ have the same index $k$ say, and for $1 \le i \le k$, $\tau(i)$ is a vertex of $\sigma(v(i))$; and for $1 \le i < j \le k$, $\tau(i) = \tau(j)$ if and only if $v(i) = v(j)$.

If $(G, \tau)$ is a rooted graph of index $k$, its underlying graft is $(G, T)$, where $T = \{\tau(1), \dots, \tau(k)\}$. It is easy to see that if $(H, v)$ is simulated in $(G, \tau)$ then $(H, U)$ is simulated in $(G, T)$, where $(H, U)$ and $(G, T)$ are the respective underlying grafts.

Again, the relation ``is simulated in'' is a quasi-order on any class of rooted graphs; and again if it is a well-quasi-order we say the class of rooted graphs is well-rooted. We need the following lemma.

\stmt
(4.3) \textit{Let $\mathscr{A}$ be a well-rooted class of grafts, and let $\mathscr{B}$ be a class of rooted graphs of bounded index, such that for each member of $\mathscr{B}$, its underlying graft is in $\mathscr{A}$. Then $\mathscr{B}$ is well-rooted.}

\textit{Proof.} It is enough to prove that for every number $k$, the class of members of $\mathscr{B}$ with index $k$ is well-rooted, since $\mathscr{B}$ is the union of only finitely many such classes.

Thus we may assume that all members of $\mathscr{B}$ have index $k$. Let $(G_1, \tau_1)$, $(G_2, \tau_2), \dots$ be a countable sequence of members of $\mathscr{B}$. For each $i$, let
